\documentclass{article}

\input{header}

\title{CSCI 2590: Assignment 4}
\author{Yvonne Yu \\ yy5919 \\ \href{https://github.com/Yvonne-Yu217/CS2590_hw_3_yy}{https://github.com/Yvonne-Yu217/CS2590\_hw\_3\_yy}}
\date{}

\colmfinalcopy
\begin{document}
\maketitle

\section*{Q0.1 Repository Link}
GitHub repository for both Part I and Part II:
\href{https://github.com/Yvonne-Yu217/CS2590_hw_3_yy}{https://github.com/Yvonne-Yu217/CS2590\_hw\_3\_yy}

\section*{Q2.1 Transformation Design}
I use a token-level stochastic transformation that combines synonym replacement with typo fallback.
For each alphabetic token, with probability $p=0.15$, I attempt to transform it. If WordNet provides
valid synonyms different from the original form, I replace the token with a randomly chosen synonym.
If no synonym is available and token length is greater than 1, I create a typo via adjacent character swap.
Non-alphabetic tokens are left unchanged.

This is a reasonable transformation because it simulates realistic user variation: paraphrasing and
small typing mistakes. In most cases sentiment should be preserved, so it serves as an OOD robustness test
without changing labels.

\section*{Q3.1 Data Augmentation Report and Analysis}
\textbf{Accuracy results}
\begin{itemize}
  \item Base model on original test: 0.93448
  \item Base model on transformed test: 0.90840
  \item Augmented model on original test: 0.93432
  \item Augmented model on transformed test: 0.92108
\end{itemize}

\textbf{Analysis}
\begin{itemize}
  \item Performance on transformed data improved after augmentation: $0.92108 - 0.90840 = +0.01268$.
  \item Performance on original data was nearly unchanged: $0.93432 - 0.93448 = -0.00016$.
  \item Intuitively, augmentation teaches the model invariance to surface perturbations while preserving
  core sentiment cues.
  \item One limitation: this augmentation is lexical/local and does not address broader distribution shifts
  such as discourse changes, long-range syntax changes, or sarcasm.
\end{itemize}

\section*{Part II. Q4 Data Statistics and Processing}
\begin{table}[h!]
\centering
\begin{tabular}{lcc}
\toprule
Statistics Name & Train & Dev \\
\midrule
Number of examples & 4225 & 466 \\
Mean sentence length (T5 tokens) & 18.10 & 18.07 \\
Mean SQL query length (T5 tokens) & 216.37 & 210.05 \\
Vocabulary size (natural language) & 792 & 466 \\
Vocabulary size (SQL) & 555 & 395 \\
\bottomrule
\end{tabular}
\caption{Data statistics before pre-processing.}
\label{tab:data_stats_before}
\end{table}

\begin{table}[h!]
\centering
\begin{tabular}{lcc}
\toprule
\textbf{T5 fine-tuned model} & Train & Dev \\
\midrule
Model name & google-t5/t5-small & google-t5/t5-small \\
Mean sentence length after prompt (T5 tokens) & 37.10 & 37.07 \\
Mean SQL query length (T5 tokens) & 216.37 & 210.05 \\
Vocabulary size (natural language, prompted) & 802 & 476 \\
Vocabulary size (SQL) & 555 & 395 \\
\bottomrule
\end{tabular}
\caption{Data statistics after pre-processing.}
\label{tab:data_stats_after}
\end{table}

\newpage

\section*{Q5 T5 Fine-tuning Details}
\begin{table}[h!]
\centering
\begin{tabular}{p{3.5cm}p{10cm}}
\toprule
Design choice & Description \\
\midrule
Data processing & Input prompt format: ``Translate the question to SQL. Question: <NL> SQL:''.
Decoder targets are ground-truth SQL sequences plus EOS token. \\
Tokenization & Default \texttt{google-t5/t5-small} tokenizer for both encoder and decoder.
Dynamic padding in collate functions. \\
Architecture & Fine-tuned pretrained T5-small end-to-end (all parameters trainable). \\
Hyperparameters & AdamW, learning rate $2\times10^{-5}$, linear scheduler, warmup epochs = 1,
batch size = 16, test batch size = 16, max epochs = 4, patience = 2. Generation used beam search
(num\_beams=4, max\_new\_tokens=256, deterministic decoding). \\
\bottomrule
\end{tabular}
\caption{Best-performing fine-tuned T5 configuration.}
\label{tab:t5_results_ft}
\end{table}

\section*{Q6 Results and Error Analysis}
\paragraph{Quantitative Results}
\begin{table}[h!]
\centering
\begin{tabular}{lcc}
  \toprule
  System & Query EM & F1 score \\
  \midrule
  \multicolumn{3}{l}{\textbf{Dev Results}} \\
  \midrule
  T5 fine-tuned (v6, best epoch) & 27.47 & 45.13 \\
  \midrule
  \multicolumn{3}{l}{\textbf{Test Results}} \\
  \midrule
  T5 fine-tuned (hidden leaderboard) & \textit{Fill from Gradescope} & \textit{Fill from Gradescope} \\
  \bottomrule
\end{tabular}
\caption{Development and test results.}
\label{tab:results}
\end{table}

\paragraph{Qualitative Error Analysis}
\begin{table}[h!]
  \centering
  \begin{tabular}{p{2cm}p{3cm}p{5cm}p{3cm}}
    \toprule
    \textbf{Error Type} & \textbf{Example} & \textbf{Description} & \textbf{Statistics} \\
    \midrule
    Syntax error & malformed SQL query & Parentheses, joins, or clause ordering errors causing execution failure. & 4/466 (best dev epoch) \\
    \midrule
    Wrong grounding & wrong city/airline/date values & Query is executable but entities/constants differ from the intent. & 289/466 (1-Record EM) \\
    \midrule
    Structural mismatch & partial or over-constrained query & Correct table family but wrong aggregation/filter structure. & 338/466 (1-SQL EM) \\
    \bottomrule
  \end{tabular}
  \caption{Representative error categories on the dev set.}
  \label{tab:qualitative}
\end{table}

\section*{Q7 Checkpoint Link}
Google Drive link containing the model checkpoint used to generate Q7 outputs:
\textbf{[Please paste your Google Drive checkpoint link here before submission]}

\section*{AI Usage}
\noindent\textbf{1. AI Tools Used}\\
ChatGPT/Copilot-style coding assistant.

\noindent\textbf{2. How AI Was Used}\\
Used for debugging environment/runtime issues, improving training/evaluation script reliability,
and polishing report language. All final code/results were validated by running experiments.

\noindent\textbf{3. Example Prompts / Interactions}\\
(1) Diagnose why SQL EM remains 0 while loss decreases. \\
(2) Propose robust evaluation-time fallback for non-SQL outputs. \\
(3) Generate reproducible overnight HPC experiment commands.

\end{document}
